\section{Limitations and open problems}
\label{sec:limitations}

The construction has four substantive limits.  Its exact palette is finite
and real: arbitrary integer states and complex phases are outside the
theorem, so the hardness source is the gate-dependent class
\(\BQPone^{G_2}\) rather than ordinary \(\BQP\), unrestricted \(\SDQCone\),
or a gate-independent perfect-completeness class.  The clean source has
\(\beta_d(\Xin)=8\); ignored mixed qubits replicate this denominator
exactly, but an intrinsically growing family of distinct initial classes
would require another exact source.  The gap is a complexity promise rather
than a practical estimate: at locality six it scales conservatively as
\(\Omega(\eta^{28}g/t^{26})\) with \(g=\Omega(L^{-3})\), and the common-copy
blow-up multiplies the vertex count by \(F=\lambda^{-2}\).  Finally, the
hardness is for the promise problem with endpoint gaps; whether our instances
also satisfy the access assumptions under which Lowe et al.\ place the
problem in \(\BQP\) is not established.

Two questions would strengthen the complexity classification.
\begin{enumerate}
\item Is normalized harmonic persistence \(\SDQCone\)-hard, as conjectured
in~\cite{lowe2026normalized}?  The missing steps are a source whose
perfect-space fraction is separated by the \(\SDQCone\) trace promise, and a
certified palette for the corresponding gate set.
\item Do the constructed filtrations satisfy an inverse-polynomial lower
bound on the nonzero principal angles between \(\cH_d(\Xin)\) and
\(\cH_d(\Xout)\), and admit efficient preparation of the initial harmonic
mixture?  Establishing both would place the instances in \(\BQP\)
by~\cite[Lemma~7]{lowe2026normalized}.
\end{enumerate}

Within its limits, the paper establishes a reusable criterion: a new clique
reduction obtains exact kernel multiplicity, an inverse-polynomial gap above
the whole kernel, and natural persistent ranks by certifying, for a finite
palette, local zero-weight kernels, private coercivity, exact filling, and
interfaces, with no growing-dimensional perturbation argument.
